\documentclass[11pt]{article}

\usepackage[margin=1in]{geometry}
\usepackage{amsmath,amssymb}
\usepackage{bm}
\usepackage{graphicx}
\usepackage{booktabs}
\usepackage{hyperref}
\usepackage{xcolor}
\hypersetup{colorlinks=true,linkcolor=blue,urlcolor=blue}

\newcommand{\vu}{\bm{u}}
\newcommand{\vf}{\bm{f}}
\newcommand{\vx}{\bm{x}}
\newcommand{\vk}{\bm{k}}
\newcommand{\vr}{\bm{r}}
\newcommand{\grad}{\nabla}
\newcommand{\dt}{\Delta t}
\newcommand{\dx}{\Delta x}

\title{\textbf{CellFlow: Technical Manual}\\[0.3em]
\large A hydrodynamically coupled agent-based cell biophysics simulator}
\author{CellFlow project}
\date{\today}

\begin{document}
\maketitle
\tableofcontents
\bigskip

\section{Overview}

CellFlow is a two-dimensional, agent-based simulator that couples discrete
cells to continuous chemical fields and to a low-Reynolds-number fluid. Each
cell is an autonomous agent that grows, divides, dies, consumes and secretes
chemicals, and interacts mechanically with its neighbours. The distinguishing
feature of CellFlow is a \emph{two-way} coupling to the fluid: cells exert
forces on the fluid, the fluid transports the chemical fields and moves the
cells, cells \emph{sense} the local fluid shear (mechanotransduction), and cells
\emph{remodel} the fluid's permeability by depositing extracellular matrix.
Cell biology also feeds back on mechanics: a cell's proliferation can be gated by
its own compressive contact pressure (contact inhibition / homeostatic pressure,
Section~\ref{sec:pressure}).

The state of the system at time $t$ consists of
\begin{itemize}
  \item a population of cells $\{c_i\}$, each with position $\vx_i$, radius
        $R_i$, accumulated nutrient $n_i$, polarity $\phi_i$, deviatoric shape
        strain $\bm{\varepsilon}_i$, cell type $\tau_i$, and phase;
  \item scalar fields on a uniform grid: nutrient $N(\vx)$, attractant
        $A(\vx)$, and extracellular matrix density $E(\vx)$;
  \item the fluid velocity field $\vu(\vx)$ and drag field $\alpha(\vx)$.
\end{itemize}

Throughout, the domain is the square $[0,L]^2$ discretized into a
$G\times G$ grid with spacing $\dx = L/G$. Time advances in steps of size $\dt$.

\section{The simulation step}\label{sec:step}

One time step executes the following operations in order
(see \texttt{CellSimulation.\_simulation\_step}):
\begin{enumerate}
  \item Assemble cell arrays (positions, radii).
  \item Compute forces: chemotactic propulsion, adhesion, repulsion, and
        post-division forces (Section~\ref{sec:forces}).
  \item Solve the fluid velocity from the force distribution
        (Section~\ref{sec:fluid}).
  \item[3b.] Mechanotransduction: update cell polarity from the fluid strain
        rate (Section~\ref{sec:mechano}).
  \item[3c.] Cell-shape mechanics: update deviatoric strain from contact stress
        (Section~\ref{sec:shape}).
  \item Check the advective CFL condition.
  \item Advect the scalar fields with the fluid velocity.
  \item Diffuse the scalar fields.
  \item Cell biology: nutrient uptake, secretion, metabolism, growth, phase,
        death (Section~\ref{sec:biology}).
  \item Move cells with the mobility relation (Section~\ref{sec:motion}).
  \item Handle division and death; resolve overlaps; enforce boundaries.
\end{enumerate}
The fluid is quasi-static (Stokes/Brinkman): it is recomputed each step from the
instantaneous force distribution, with no time history.

\section{Cell biology}\label{sec:biology}

\subsection{Nutrient and growth}
Each cell consumes nutrient over its circular footprint at a per-cell rate
$k_i \sim \mathcal{N}(0.2, 0.05)$, accumulating internal nutrient $n_i$. Per step
the cell takes up $\Delta n_i^{+}$ (the integrated uptake over its area, capped
by the locally available nutrient), pays a basal metabolic cost
$b\,\dt$ with $b=0.02$, and optionally secretes attractant at rate $s=1.0$:
\begin{equation}
  n_i \leftarrow n_i + \Delta n_i^{+} - b\,\dt .
\end{equation}
A cell dies (is removed) when $n_i < 0$.

\paragraph{Uptake kinetics.} The local uptake rate as a function of the ambient
concentration $c$ is selectable. By default it is \emph{first-order} (linear),
$q(c)=k_i\,c$. Setting a half-saturation constant $K_m$ (per cell, or globally
via \texttt{nutrient\_uptake\_saturation}) switches to \emph{Michaelis--Menten /
Monod} saturating kinetics,
\begin{equation}
  q(c) = k_i\,c\,\frac{K_m}{K_m + c},
\end{equation}
so that $k_i$ is the low-concentration rate constant ($q\to k_i c$ as
$c\ll K_m$) and the uptake saturates at $k_i K_m$ for $c\gg K_m$ (a cell cannot
consume arbitrarily fast). The linear law is recovered exactly as
$K_m\to\infty$. Saturating uptake reshapes the nutrient profile and is the
realistic regime for diffusion-limited aggregates and colony fronts.

\subsection{Size models}
CellFlow provides two growth laws, selected by \texttt{growth\_model}.

\paragraph{Linear (legacy).} The radius is linear in accumulated nutrient,
\begin{equation}
  R_i = \mathrm{clip}\!\left(R_{\min} + n_i\,\frac{R_{\max}-R_{\min}}{100},\,
  R_{\min},\, R_{\max}\right),
\end{equation}
with $R_{\min}=2$, $R_{\max}=4$. Because area $\propto R^2$, area grows
quadratically in $n_i$, and division (below) is \emph{not} area-conserving.

\paragraph{Area-conserving (recommended for physical fidelity).} The cell
\emph{area} grows linearly with accumulated nutrient (mass), and the radius is
recovered from the area:
\begin{equation}
  A_i = \pi R_{\max}^2 \,\frac{n_i}{100}, \qquad
  R_i = \min\!\left(\sqrt{A_i/\pi},\, R_{\max}\right).
\end{equation}
At the division threshold ($n_i=100$, $R_i=R_{\max}$), halving $n_i$ halves the
area, so each of the two daughters has area $A_{\max}/2$ and radius
$R_{\max}/\sqrt{2}$; the two daughters therefore conserve the mother's area
exactly.

\subsection{Division and death}
A cell enters the \texttt{DIVISION} phase when $R_i \ge R_{\max}$. On division it
halves its accumulated nutrient, returns to \texttt{GROWTH}, and creates a
daughter placed at contact distance $R_{\text{parent}} + R_{\text{daughter}}$
along a random unit direction (so the daughter is born just touching, not
overlapping). The daughter inherits cell type, polarity, and growth mode; its
deviatoric shape strain starts at zero (a newborn is round). For a brief window
after division the mother--daughter pair feel a mutual separating force.

\subsection{Mechanical feedback on growth}\label{sec:pressure}
By default a cell's biology is mechanically blind: it exerts and feels contact
forces but proliferates regardless of how crowded it is. Optionally
(\texttt{enable\_pressure\_inhibition}) cells respond to compressive stress via
\emph{contact inhibition of proliferation} / a homeostatic pressure. Each cell is
assigned a scalar (virial) contact pressure from its repulsive overlaps,
\begin{equation}
  P_i = \frac{1}{2 A_i}\sum_{j}\,\vf_{ij}\cdot\vr_{ij}, \qquad A_i = \pi R_i^2,
\end{equation}
where the sum runs over overlapping neighbours and $\vf_{ij}\cdot\vr_{ij}\ge 0$
under repulsion (compression). Unlike the net force, $P_i$ does not cancel under
isotropic squeezing, so it measures mechanical crowding. When $P_i$ exceeds
\texttt{pressure\_threshold} the cell grows to full size but does not divide
(its stored nutrient is capped at the division threshold). The compressed colony
core therefore arrests while the uncrowded rim keeps proliferating, and the
population settles to a homeostatic density rather than overgrowing.

\section{Chemical fields}\label{sec:fields}

Each scalar field $c \in \{N, A\}$ obeys advection--diffusion with cellular
source/sink terms,
\begin{equation}
  \frac{\partial c}{\partial t} + \vu\cdot\grad c
  = D_c \grad^2 c + S_c(\vx,t),
\end{equation}
solved by operator splitting each step: advection, then diffusion, then the
cellular uptake/secretion. Boundary conditions for the nutrient field are
selectable: \emph{Neumann} (no-flux, isolated system) or \emph{Dirichlet}
(edges held at a constant reservoir value, modelling an external medium).

\paragraph{Diffusion.} Two solvers are available, selected by
\texttt{diffusion\_solver}. The default \emph{explicit} five-point (FTCS) scheme,
\begin{equation}
  c^{m+1}_{ij} = c^{m}_{ij}
  + \frac{D_c \dt}{\dx^2}\left(c_{i-1,j}+c_{i+1,j}+c_{i,j-1}+c_{i,j+1}-4c_{ij}\right),
\end{equation}
clamped to be non-negative, is stable only when the diffusion number
$D_c\dt/\dx^2 \le 1/4$ (the code warns otherwise).

The \emph{implicit} solver (\texttt{diffusion\_solver='implicit'}) is an
alternating-direction implicit (ADI, Peaceman--Rachford) scheme: each step is two
half-steps, the first implicit in $x$ and explicit in $y$, the second the reverse,
\begin{align}
  \big(\mathcal{I} - \tfrac{r}{2}\,\delta_x^2\big)c^{*}
   &= \big(\mathcal{I} + \tfrac{r}{2}\,\delta_y^2\big)c^{m}, \\
  \big(\mathcal{I} - \tfrac{r}{2}\,\delta_y^2\big)c^{m+1}
   &= \big(\mathcal{I} + \tfrac{r}{2}\,\delta_x^2\big)c^{*},
\end{align}
with $r = D_c\dt/\dx^2$ and $\delta^2$ the second-difference operator. Each
half-step is a set of independent tridiagonal systems solved by the Thomas
algorithm ($O(M)$), parallelized over lines. ADI is \emph{unconditionally stable}
(any $\dt$), second order in space and time, and preserves the discrete steady
state exactly (a uniform field at the Dirichlet value is a fixed point), so it
reaches the quasi-steady profile in far fewer steps than the explicit scheme and
without a CFL limit on $D$. Because CellFlow advects fields by the flow, the
overall scheme is IMEX (implicit diffusion combined with the semi-Lagrangian
advection below), which keeps the cell--fluid coupling intact. ADI is
$L^2$-stable rather than strictly maximum-principle-preserving, so very large
steps can produce small ($\sim$few-percent) over/undershoots, and with a fixed
step it converges to the exact steady state gradually (no parameter cycling).

Under no-flux (Neumann) boundaries the explicit scheme is conservative in the
interior (verified) but the node-centred ghost treatment is not strictly
flux-conservative at the walls.

\paragraph{Advection.} A semi-Lagrangian scheme back-traces each grid point
along $-\vu\,\dt$ and interpolates bilinearly. It is unconditionally stable but
diffusive and first-order; it conserves mass in the interior under smooth flow
(verified to $\sim 10^{-5}$) but loses mass where material is advected into the
clamped domain boundary.

\section{Cell forces}\label{sec:forces}

Forces on the cells are separated into a \emph{propulsion} component (active,
chemotactic) and a \emph{monopolar} component (passive, mechanical), so that
the dipole hydrodynamic model can treat propulsion as force-free.

\paragraph{Chemotactic propulsion.} Each cell samples the area-averaged nutrient
gradient and combines it with a random walk to set a direction, then propels at
fixed magnitude $F_{\text{prop}}$:
\begin{equation}
  \hat{\vf}_i \propto \chi\,\overline{\grad N}_i + w\,\bm{\eta}_i, \qquad
  \vf^{\text{prop}}_i = F_{\text{prop}}\,\frac{\hat{\vf}_i}{\lvert\hat{\vf}_i\rvert},
\end{equation}
where $\chi$ is the chemotactic coefficient, $w$ the random-walk speed, and
$\bm{\eta}_i$ standard normal noise. The noise is drawn once per step in the main
thread (seedable) and passed into the parallel kernel, so results are
reproducible and independent of thread count.

\paragraph{Repulsion (volume exclusion).} For overlapping cells
($d_{ij} < R_i + R_j$), with overlap $\delta_{ij}=R_i+R_j-d_{ij}$,
\begin{equation}
  \vf^{\text{rep}}_i = -\sum_{j} k_{\text{rep}}\,
  \exp\!\left(\frac{3\delta_{ij}}{R_i+R_j}\right)\hat{\vr}_{ij},
\end{equation}
directed away from the neighbour.

\paragraph{Adhesion.} For cells within an outer band
$R_i+R_j < d_{ij} < \lambda(R_i+R_j)$ (cutoff factor $\lambda$), an attractive
force linear in the separation beyond contact,
\begin{equation}
  \vf^{\text{adh}}_i = \sum_j k_{\text{adh}}\,(d_{ij}-(R_i+R_j))\,\hat{\vr}_{ij}.
\end{equation}

\paragraph{JKR adhesive contact (\texttt{contact\_model='jkr'}).} The
repulsion/adhesion pair above has no derivation and a discontinuity at contact:
repulsion switches on at $d_{ij}=R_i+R_j$ with magnitude $k_{\text{rep}}$, while
adhesion acts only \emph{outside} contact. Measured on a packed colony the mean
repulsive force is $56$ against a mean adhesive force of $0.11$, a factor of
$500$, and the resulting cohesive well is $0.75$ deep. Johnson--Kendall--Roberts
contact replaces both with a single continuous law, parameterised by the contact
radius $a$:
\begin{align}
  F(a)      &= \frac{4E^\ast a^3}{3\hat R} - \sqrt{8\pi w E^\ast a^3},\\
  \delta(a) &= \frac{a^2}{\hat R} - \sqrt{\frac{2\pi w a}{E^\ast}},
\end{align}
with $\hat R = R_iR_j/(R_i+R_j)$ the reduced radius, $E^\ast$ the reduced modulus
and $w$ the work of adhesion. Given the overlap $\delta$, $a$ is obtained on the
stable branch $a\ge a_{\min}$ by safeguarded Newton iteration.

Two features distinguish it from the law it replaces. First, a genuine cohesive
well. Second, $\delta(a)$ is non-monotonic, with a minimum at
$a_{\min}=(cR/4)^{2/3}$, $c=\sqrt{2\pi w/E^\ast}$: contact therefore persists to
\emph{negative} overlap, so cells resist being pulled apart. The two closed-form
landmarks, both reproduced to six significant figures, are
\begin{equation}
  F_c = -\tfrac{3}{2}\pi w \hat R \quad\text{(pull-off, load control)},\qquad
  \delta_c = -\tfrac{3}{4}\left(\frac{\pi^2w^2\hat R}{E^{\ast 2}}\right)^{1/3}
  \quad\text{(snap, displacement control)} .
\end{equation}
These are \emph{different} points; since the simulation is displacement
controlled, the neck snaps at $\delta_c$, where the force is not $F_c$. Setting
$w=0$ recovers Hertz exactly. JKR carries its own repulsion continuously, so the
overlap projection is unnecessary---and harmful, since it overrides the force
balance---and defaults off.

\paragraph{Timestep bound.} JKR contacts are stiff. With an overdamped velocity
law, a contact of stiffness $k=\mathrm{d}F/\mathrm{d}\delta$ relaxes at rate
$k/\gamma$, so explicit Euler requires
\begin{equation}
  \dt < \frac{2\gamma}{k_{\max}},\qquad
  k_{\max}=\max_{\text{contacts}}\left|\frac{\mathrm{d}F}{\mathrm{d}\delta}\right| .
\end{equation}
This matters because violating it does \emph{not} blow up loudly: the pack
fragments quietly while still producing plausible shape statistics. The bound is
computed by \texttt{jkr\_stable\_dt} and the simulation warns once when exceeded.
Note $\mathrm{d}\delta/\mathrm{d}a=0$ at the snap point, so the stiffness diverges
there; a pack built at zero overlap rather than at the JKR equilibrium overlap
starts every contact near that singularity, and the admissible timestep collapses
by an order of magnitude.

\paragraph{Differential adhesion.} With cell types and a symmetric adhesion
matrix $\bm{D}$, the scalar strength $k_{\text{adh}}$ is replaced by
$D_{\tau_i\tau_j}$. Following the Differential Adhesion Hypothesis, the
interfacial tension between types is
$\sigma_{ab} = D_{ab} - \tfrac{1}{2}(D_{aa}+D_{bb})$, which drives cell sorting
and engulfment. A uniform matrix reproduces the scalar law exactly.

\paragraph{Neighbour lists.} All pairwise force loops use a uniform linked-cell
(bin) structure when enabled (default), reducing the cost from $O(N^2)$ to
$O(N)$; the binned forces are identical to the brute-force result to
$\sim 10^{-10}$.

\paragraph{Overlap resolution.} After the force-driven move, residual overlaps
are removed by a geometric projection: each overlapping pair is pushed apart by
half the overlap, iterated a configurable number of times per step. With
neighbour lists enabled (default) this uses a parallel cell-list \emph{Jacobi}
sweep---each cell accumulates the half-overlap push from its overlapping
neighbours (read-only on positions) and all displacements are applied at
once---which is $O(N)$ rather than the $O(N^2)$ sequential (Gauss--Seidel)
sweep, and preserves the centre of mass exactly (equal-and-opposite pushes);
an isolated pair is resolved in a single sweep. The degenerate
exactly-coincident case is separated along a deterministic golden-angle
direction, preserving reproducibility.

\section{Hydrodynamics}\label{sec:fluid}

The fluid is governed by quasi-static, incompressible low-Reynolds flow. Three
solvers are available, selected by \texttt{fluid\_model} and related options.

\subsection{Free-space regularized Stokeslet (legacy)}
The velocity is the sum of 2D regularized Stokeslets (Cortez) over the cell
forces. For a force $\vf$ at $\vx_0$, with $\vr=\vx-\vx_0$,
$r_\epsilon^2 = \lvert\vr\rvert^2+\epsilon^2$ and regularization
$\epsilon = 2\dx$,
\begin{equation}
  \vu(\vx) = \frac{1}{4\pi\mu}\Big[
  -\vf\big(\ln r_\epsilon - \tfrac{\epsilon^2}{r_\epsilon^2}\big)
  + (\vf\cdot\vr)\,\frac{\vr}{r_\epsilon^2}\Big].
\end{equation}
A force-dipole variant represents active propulsion as two equal-and-opposite
Stokeslets (no net force on the fluid). \textbf{Caveat:} in 2D the Stokeslet
velocity grows like $\ln r$ for a net force (the Stokes paradox), so this model
has no far-field decay and no confinement; it is retained for backward
compatibility and force-free (dipole) studies.

\subsection{FFT Brinkman solver (recommended)}
The preferred model solves the incompressible Brinkman equations on a periodic
grid,
\begin{equation}
  -\mu\grad^2\vu + \alpha\,\vu + \grad p = \vf, \qquad \grad\cdot\vu = 0,
  \label{eq:brinkman}
\end{equation}
where $\mu$ is the viscosity and $\alpha$ the substrate drag. The drag
introduces a screening length $\delta = \sqrt{\mu/\alpha}$ beyond which
hydrodynamic interactions decay; this regularizes the 2D paradox. In Fourier
space, for each wavevector $\vk\neq 0$,
\begin{equation}
  \hat{\vu}(\vk) = \frac{\mathsf{P}(\vk)\,\hat{\vf}(\vk)}{\mu\lvert\vk\rvert^2+\alpha},
  \qquad
  \mathsf{P}(\vk) = \mathsf{I} - \frac{\vk\vk^{\mathsf{T}}}{\lvert\vk\rvert^2},
  \label{eq:brinkman-fft}
\end{equation}
where $\mathsf{P}$ is the Leray (divergence-free) projection. The mean mode is
$\hat{\vu}(0)=\hat{\vf}(0)/\alpha$ (finite for $\alpha>0$). The cost is
$O(M\log M)$ for $M=G^2$ grid points. The Nyquist wavenumber is zeroed on even
grids so the projection is consistent (otherwise the velocity is not
divergence-free). The solver is verified against an analytic manufactured
solution to machine precision and is divergence-free to $\sim 10^{-9}$.

\subsection{Free-slip box (non-periodic walls)}\label{sec:freeslip}
Setting \texttt{fluid\_boundary='freeslip\_box'} replaces the periodic boundaries
with free-slip (no-penetration, stress-free) walls on all four sides of the
square domain,
\begin{equation}
  \vu\cdot\bm{n} = 0, \qquad
  \big(\mathsf{\sigma}\bm{n}\big)\cdot\bm{t} = 0
  \quad\text{on each wall.}
\end{equation}
A free-slip wall is a plane of mirror symmetry, so the box solution is the
restriction of a periodic solution on the doubled domain $[0,2L]^2$ whose force
is extended with the parity $f_x$ odd across the $x$-walls and even across the
$y$-walls, and $f_y$ even across the $x$-walls and odd across the $y$-walls (the
extensions underlying the discrete sine/cosine transforms). The Brinkman
multiplier in \eqref{eq:brinkman-fft} and the Leray projection $\mathsf{P}$ both
map this symmetry subspace to itself, so the solved velocity inherits the
symmetry---$u_x$ odd in $x$ and even in $y$, $u_y$ even in $x$ and odd in
$y$---which is exactly free slip, and remains divergence-free. Reusing the
verified periodic solver on the mirror-extended grid keeps the cost
$O(M\log M)$ (a constant factor of $4$). This path is verified against a
manufactured free-slip mode to machine precision. It is currently restricted to
constant drag $\alpha$ (not combined with the poroelastic ECM solver), and no-slip
walls would instead require a Chebyshev/influence-matrix formulation.

\subsection{Variable-coefficient (poroelastic) Brinkman}
When cells remodel the matrix (Section~\ref{sec:ecm}), the drag becomes a field
$\alpha(\vx)$, and \eqref{eq:brinkman} is no longer diagonal in Fourier. Writing
$\alpha(\vx) = \alpha_0 + d(\vx)$ with $\alpha_0 = \max_{\vx} \alpha(\vx)$ and
$d \le 0$, the solution satisfies the fixed point
\begin{equation}
  \vu = \mathsf{A}_0^{-1}\mathsf{P}\big(\vf - d\,\vu\big),
\end{equation}
where $\mathsf{A}_0^{-1}\mathsf{P}$ is precisely the constant-$\alpha_0$ FFT
solve \eqref{eq:brinkman-fft}. This Picard iteration contracts because
$\lvert d\rvert/(\mu\lvert\vk\rvert^2+\alpha_0) < 1$, with rate
$\sim 1 - \alpha_{\min}/\alpha_{\max}$. It is efficient for moderate drag
contrast; very high contrast would warrant a preconditioned Krylov solver
(future work). It is verified by constant-$\alpha$ consistency, a
spatially-varying manufactured solution (to $10^{-6}$), and flow screening
through a high-drag patch.

\subsection{Growth-driven expansion (volumetric source)}\label{sec:growthsource}
Cells that grow create material in place. Setting
\texttt{enable\_growth\_source} adds this to the fluid as a volumetric source,
replacing the incompressibility constraint in \eqref{eq:brinkman} with
\begin{equation}
  -\mu\grad^2\vu + \alpha\,\vu + \grad p = \vf, \qquad
  \grad\cdot\vu = s(\vx),
  \label{eq:brinkman-source}
\end{equation}
where $s$ is the rate at which cells produce area per unit area. The colony then
displaces its surroundings through a long-range \emph{pressure} field rather than
only through the local steric overlap projection---the Darcy/Saffman--Taylor
mechanism used by continuum colony models, $\bm{v} = -K_p\grad p$. Note that
\eqref{eq:brinkman-source} reduces to Darcy's law when $\delta$ is small compared
with the colony, so this is the continuum mechanism itself and not an analogy.

Decompose $\vu = \bm{w} + \grad\phi$ with $\grad^2\phi = s$ and
$\grad\cdot\bm{w} = 0$. The potential part contributes
$-\mu\grad^2\grad\phi + \alpha\grad\phi = -\mu\grad s + \alpha\grad\phi$. For
\emph{constant} $\alpha$ both terms are gradients and are absorbed into the
pressure, so the two parts separate exactly and
\begin{equation}
  \hat{\vu}(\vk) = \frac{\mathsf{P}(\vk)\hat{\vf}(\vk)}{\mu\lvert\vk\rvert^2+\alpha}
  \;-\; \frac{i\vk\,\hat{s}(\vk)}{\lvert\vk\rvert^2}.
\end{equation}
For variable $\alpha(\vx)$ the term $\alpha\grad\phi$ is not a gradient; it
survives as a body force, so the split remains exact provided the solenoidal part
is driven by $\vf - \alpha(\vx)\grad\phi$. Both cases are implemented.

\paragraph{Building $s$ from cell growth.} A cell whose area goes from $A_k^{n}$
to $A_k^{n+1}$ over $\dt$ produces area at rate $(A_k^{n+1}-A_k^{n})/\dt$, spread
over its footprint, so that
$\int s\,\mathrm{d}A = \sum_k \mathrm{d}A_k/\mathrm{d}t$. Shrinking cells
contribute a sink, which is correct. This requires a deposition kernel that
conserves the total: the general-purpose secretion kernel normalizes by the
\emph{analytic} point count $\pi(R/\dx)^2$ rather than the number of grid points
it actually writes to, which is harmless for a secretion amplitude but misstates
the deposited total by tens of percent at low resolution---and here the integral
of $s$ \emph{is} the volume the colony pushes out, so an amplitude error is an
error in the driving flow (Gauss). A separate conserving kernel is used.

Because $\grad\cdot\vu = s$ has a periodic solution only if $s$ has zero mean,
the mean is subtracted: physically, a uniform far-field drain removing the fluid
the colony displaces. This is accurate while the colony is small compared with the
box. The source is built from the area produced during a step and applied at the
next one (a one-step lag; the Stokes solve is quasi-static, so this is a splitting
choice rather than an accumulating error).

Verification: a manufactured potential is recovered to machine precision; the
divergence matches the mean-removed source (exactly for a band-limited source, and
to the source's Nyquist content otherwise); the flux through a circle equals the
enclosed source (Gauss); and the far field follows the 2D point-source law
$u_r = Q/2\pi r$ to within $3\%$.

\emph{Double counting.} With this enabled, growth expands the colony both through
the pressure field and through the overlap projection of
Section~\ref{sec:forces}. Reduce \texttt{overlap\_iterations}, or scale
\texttt{growth\_source\_strength}, so the expansion is not counted twice; the
appropriate split depends on how much of the tissue mechanics the fluid should
carry.

\subsection{Immersed Boundary coupling}\label{sec:ibm}
Cells couple to the grid fluid by the Immersed Boundary Method. A cell force
$\vf_i$ is \emph{spread} to a grid force density with a per-cell Gaussian blob of
\emph{physical} width $\sigma_i = \zeta R_i$ (where $\zeta$ is
\texttt{ibm\_reg\_factor}):
\begin{equation}
  \vf(\vx) = \sum_i \vf_i\,\varphi_{\sigma_i}(\vx-\vx_i),
\end{equation}
normalized per cell so that total force is conserved exactly. The solved
velocity is \emph{interpolated} back to each cell with the same kernel,
\begin{equation}
  \vu_i = \sum_{\vx} \vu(\vx)\,\hat{\varphi}_{\sigma_i}(\vx-\vx_i),
\end{equation}
with weights summing to one (partition of unity), so a uniform field is
recovered exactly. Tying the regularization to the cell radius (rather than the
grid spacing, as a grid-width kernel would) makes the single-cell self-mobility
\emph{grid convergent}: it approaches a finite limit as $\dx\to 0$, whereas a
grid-tied kernel drifts (logarithmically). The same field $\vu$ both advects the
scalar fields and moves the cells, removing any kinematic inconsistency.

\subsection{Cell motion}\label{sec:motion}
In the Brinkman/IBM path, cell velocities are the interpolated fluid velocity at
the cell centres. In the legacy Stokeslet path, the velocity is a mobility
relation combining a self term (Stokes drag) and hydrodynamic interactions:
\begin{equation}
  \vu_k = \frac{\vf_k}{6\pi\mu R_k}
  + \sum_{j\neq k} \mathsf{G}_{2D}(\vx_k,\vx_j)\,\vf_j .
\end{equation}
Cells advance by forward Euler, $\vx_k \leftarrow \vx_k + \vu_k\,\dt$, which is
first order in $\dt$ (verified).

\paragraph{Local friction law (\texttt{velocity\_model='friction'}).} Both paths
above are \emph{non-local}: a cell's velocity is set by a field, so neighbouring
cells share it. Under the Brinkman/IBM path this is exact---$\vu_k$ is the
interpolated fluid velocity---and it has a consequence worth stating plainly:
two neighbours cannot acquire relative velocity at the cell scale, and therefore
can never exchange places. The centre-based-model standard instead solves a local
friction balance,
\begin{equation}
  \gamma_{\text{sub}}\,\vu_i
  + \sum_j \gamma_{\text{cc}}\,w_{ij}\,(\vu_i-\vu_j) = \vf_i ,
\end{equation}
with $\gamma_{\text{sub}}$ substrate drag, $\gamma_{\text{cc}}$ cell--cell
friction and $w_{ij}$ a contact weight tapering to zero at the cutoff. The
operator is $\gamma_{\text{sub}}\mathsf{I}+\mathsf{L}$ with $\mathsf{L}$ the
weighted graph Laplacian of the contact network---symmetric positive definite for
$\gamma_{\text{sub}}>0$---and is solved matrix-free by Jacobi-preconditioned
conjugate gradient. Galilean invariance is structural: $\mathsf{L}$ annihilates
uniform velocity, so only substrate drag breaks translation invariance. An
isolated cell recovers $\vu=\vf/\gamma_{\text{sub}}$ exactly.

\paragraph{The fluid path is a low-pass filter on force.} This is a property of
the Brinkman solver, not a bug, but it governs which physics can be represented.
The transfer function (Section~\ref{sec:fluid}) gives
\begin{equation}
  \frac{\lvert\vu(k)\rvert}{\lvert\vu(k\!\to\!0)\rvert}
  = \frac{1}{1+(k\delta)^2},
\end{equation}
so a force varying at wavenumber $k$ produces velocity suppressed by
$1+(k\delta)^2$. At $\mu=500$, $\delta=14$ and a cell spacing of $4.32$ that is a
factor of $416$ at the cell scale. Measured consequence: a random propulsive
force of $150$ per cell---three times the mean repulsive force---changes the mean
displacement over $500$ steps from $3.154$ to $3.159$, i.e. not at all. Cell-scale
rearrangement, and hence surface-tension-driven shape relaxation, cannot occur
through this velocity law at any screening length. The friction law has no such
filter. Choose the fluid path for suspensions and bioreactors, where momentum
genuinely travels through the medium, and the friction path for dense tissue on a
substrate.

\section{Interfacial surface tension}\label{sec:surften}

The model has no curvature-dependent term of its own. Every mechanical term
scales with the expansion rate, with the measured consequence that the front
dispersion relation factorises as
$\lambda(k)=(\dot R/R)\,f(k)$, so the marginal mode $k_0=\lambda_0/c$ is a pure
number that no growth-side parameter can move (Section~\ref{sec:limits}).
Setting \texttt{surface\_tension} $>0$ supplies Young--Laplace explicitly at the
boundary, which is the closure continuum colony models use,
$p = p_0 - \sigma\kappa$.

\paragraph{Why not a grid continuum surface force.} The standard construction
(Brackbill CSF) smooths an occupancy field $\phi$, forms
$\kappa=-\nabla\!\cdot\!(\nabla\phi/\lvert\nabla\phi\rvert)$ and applies
$\sigma\kappa\bm{n}\lvert\nabla\phi\rvert$. This was implemented and rejected on
measurement: curvature is a second derivative, the cell-scale roughness of a
packed boundary dominates it, and $\kappa R$ for a disc came out between $0.07$
and $1.33$ depending on colony size and smoothing width, with no setting accurate
across sizes.

\paragraph{Band-limited boundary parameterisation.} Instead the boundary is
written as $R(\theta)$ about the colony centroid and truncated to angular modes
$k\le k_{\max}$, after which curvature follows in closed form,
\begin{equation}
  \kappa = \frac{R^2 + 2R'^2 - RR''}{\left(R^2+R'^2\right)^{3/2}},
\end{equation}
with the derivatives taken spectrally, so the result is \emph{exact} for the
truncated shape rather than a finite difference of a noisy one. Truncation is the
physically appropriate regularisation: front modes of interest are $k\lesssim 20$,
and cell-scale modes are not interface shape. Verified: $\kappa R = 1.000000$ for
discs at every radius tested, and the analytic curvature of a lobed front to
eight significant figures including the negative values in its valleys.

The force $\sigma\kappa$ is applied inward along the local normal to boundary
cells only, so it is a surface rather than a body term. For a disc it reduces to
a uniform inward $\sigma/R$, reproduced to $10^{-14}$.

\paragraph{Limitation.} $R(\theta)$ requires a colony that is star-shaped about
its centroid. That holds for lobed discs through deep-necked fingers, but would
fail for an overhanging or fragmented shape, and callers should check.

\paragraph{Scope.} With this term active an interfacial instability is
\emph{imposed} rather than emergent from cell-scale rules---as it is in the
continuum models being compared against. Section~\ref{sec:kb} gives the
measurement that distinguishes the two.

\section{Mechanotransduction}\label{sec:mechano}

Cells sense the local fluid shear and reorient. From the solved velocity, the
velocity-gradient tensor $\grad\vu$ is computed by periodic central differences
and area-averaged over each cell. The strain-rate tensor is
$\mathsf{E} = \tfrac{1}{2}(\grad\vu + \grad\vu^{\mathsf{T}})$, from which we
form the scalar shear rate and the principal (extensional) axis,
\begin{align}
  \dot\gamma &= \sqrt{2\,\mathsf{E}\!:\!\mathsf{E}}
  = \sqrt{2\big(E_{xx}^2+E_{yy}^2+2E_{xy}^2\big)}, \\
  \psi &= \tfrac{1}{2}\,\mathrm{atan2}\big(2E_{xy},\, E_{xx}-E_{yy}\big).
\end{align}
The cell polarity $\phi$ (a nematic director, defined modulo $\pi$) aligns toward
the strain axis at a rate proportional to the shear,
\begin{equation}
  \frac{d\phi}{dt} = -k_\phi\,\dot\gamma\,\sin\!\big(2(\phi-\psi)\big),
\end{equation}
with rate constant $k_\phi$ (\texttt{shear\_alignment\_rate}). The stable fixed
point is $\phi=\psi$. Optionally the polarity drives an active propulsion of
magnitude \texttt{polarity\_propulsion\_force} (rheotaxis). Verification: under
analytic simple shear $u=(\dot\gamma\,y,0)$ the sensed shear rate equals
$\dot\gamma$, the axis is $45^\circ$, and the polarity converges to $45^\circ$;
zero flow leaves the polarity unchanged.

\section{Cell-shape mechanics}\label{sec:shape}

Cells deform into area-conserving ellipses under contact pressure, while the
mechanical contacts remain circular (so the cost stays low). For each cell the
repulsive contact (Cauchy) stress is
\begin{equation}
  \bm{\sigma}_i = \frac{1}{A_i}\sum_j \vf_{ij}\otimes\vr_{ij},
\end{equation}
with deviatoric part
$\bm{s}_i = \big(\tfrac{1}{2}(\sigma_{xx}-\sigma_{yy}),\, \sigma_{xy}\big)$. The
deviatoric strain $\bm{\varepsilon}_i=(\varepsilon_{xx},\varepsilon_{xy})$
relaxes by a linear viscoelastic law toward $\chi\,\bm{s}_i$,
\begin{equation}
  \bm{\varepsilon}_i \leftarrow \bm{\varepsilon}_i
  + \big(\chi\,\bm{s}_i - \bm{\varepsilon}_i\big)\frac{\dt}{\tau},
\end{equation}
capped at a maximum aspect ratio, with compliance $\chi$
(\texttt{shape\_compliance}) and relaxation time $\tau$
(\texttt{shape\_relaxation\_time}). The ellipse semi-axes follow from the strain
magnitude $m=\lvert\bm{\varepsilon}_i\rvert$,
\begin{equation}
  a = R_i\,e^{m}, \quad b = R_i\,e^{-m}, \quad
  \theta = \tfrac{1}{2}\mathrm{atan2}(\varepsilon_{xy},\varepsilon_{xx}),
\end{equation}
so that $ab=R_i^2$ (area conserved). Verification: uniaxial compression
elongates perpendicular to the squeeze and reaches the elastic steady state
$\bm{\varepsilon}=\chi\bm{s}$; \emph{isotropic} (hydrostatic) pressure produces
\emph{no} elongation; the shape relaxes back to round when unloaded.

\section{Dynamic extracellular matrix}\label{sec:ecm}

Cells deposit ECM into a grid field $E(\vx)$ over their area; the field decays
and is mapped to the local Brinkman drag,
\begin{equation}
  E \leftarrow (1 - \kappa_d\,\dt)\,E + \text{(cell secretion)}, \qquad
  \alpha(\vx) = \alpha_{\text{base}} + \kappa_\alpha\,E(\vx),
\end{equation}
with secretion rate, decay rate $\kappa_d$, and drag coefficient $\kappa_\alpha$
set by \texttt{ecm\_secretion\_rate}, \texttt{ecm\_decay\_rate}, and
\texttt{ecm\_drag\_coeff}. The resulting variable drag field is passed to the
poroelastic solver, so cell-deposited matrix lowers the local permeability and
reroutes/screens the flow---the defining cell$\to$fluid feedback (relevant to
tumour stroma, interstitial flow, and durotaxis).

\section{Numerical methods and stability}

\begin{itemize}
  \item \textbf{Time integration:} forward Euler for cell positions and the
        ODE-like polarity/shape relaxations (first order in $\dt$).
  \item \textbf{Diffusion:} selectable --- explicit FTCS (stable for
        $D\dt/\dx^2 \le 1/4$) or implicit ADI (Peaceman--Rachford, Thomas solves,
        unconditionally stable, second order, IMEX with the flow advection).
  \item \textbf{Uptake:} first-order (linear) or Michaelis--Menten/Monod
        saturating kinetics.
  \item \textbf{Advection:} semi-Lagrangian; the code warns when the CFL number
        $\max\lvert\vu\rvert\,\dt/\dx > 1$.
  \item \textbf{Fluid:} spectral (FFT) for Brinkman, $O(M\log M)$; Picard
        iteration for variable drag.
  \item \textbf{Pairwise forces:} linked-cell neighbour lists, $O(N)$; an optional
        per-cell virial contact pressure (same linked-cell structure) drives
        mechanical feedback on proliferation.
  \item \textbf{Multi-timescale splitting:} optional decoupling of fast transport
        from the slow/expensive steps. \texttt{diffusion\_substeps} $M$ applies the
        diffusion solver $M$ times per step with $\dt/M$ (keeps the explicit scheme
        stable when $D\dt/\dx^2>\tfrac14$ overall, and relaxes the field toward
        quasi-steady). \texttt{fluid\_update\_interval} $K$ recomputes the
        expensive Brinkman solve only every $K$ steps, re-interpolating cell
        velocities from the cached field each step (bounded staleness; always
        recomputed when the ECM drag field is active).
\end{itemize}

\section{Colony-front analysis}\label{sec:front}

The module \texttt{cellflow.analysis.front} provides quantitative diagnostics for
a growing colony boundary. It takes plain arrays (positions, radii) rather than a
\texttt{CellSimulation}, so it can be tested against analytic shapes.

\subsection{Extracting the front}
The front $R(\theta)$ is sampled on $n_\theta$ equal angular bins about a given
centre, taking the \emph{outermost} cell in each bin (empty bins are filled by
periodic interpolation). Taken naively this measure is dangerous: whenever cells
detach from the colony---which happens for weak adhesion and a starving rim---a
handful of strays per bin inflates the radius, and the front reads as rough while
the actual colony is smooth. All front quantities are therefore computed on the
\textbf{largest connected cluster}, where two cells are connected when
\begin{equation}
  \lvert \vx_i - \vx_j\rvert < \eta\,(R_i + R_j),
\end{equation}
with $\eta=1.3$ by default ($\eta=1$ is exactly touching). Components are labelled
by union--find over the existing linked-cell grid, so the cost is $O(N)$.

\subsection{Mode amplitudes and growth rates}
The dimensionless amplitude of angular mode $k$ is
\begin{equation}
  a_k = \frac{2\,\lvert \mathcal{F}_k[R - \langle R\rangle]\rvert}
             {n_\theta\,\langle R\rangle},
\end{equation}
normalized so that a front $R=R_0(1+\varepsilon\cos k\theta)$ gives exactly
$a_k=\varepsilon$, independent of $R_0$. A linear instability grows as
$a_k(t)=a_k(0)e^{\lambda t}$, so $\lambda$ is obtained by least squares on
$\log a_k$ against $t$. Fitting the logarithm (rather than taking an
end-to-start ratio) is what separates genuine exponential growth from a transient
excursion: a non-exponential series shows up as a poor $r^2$ even when the
endpoints happen to lie above the start.

Two conventions coexist. The \emph{relative} rate uses $a_k=\delta_k/R$ and is the
fingering criterion (the shape deviating more and more); the \emph{absolute} rate
uses $\delta_k$. On a colony expanding at rate $\mathrm{d}\ln R/\mathrm{d}t$ the
two differ by exactly that rate.

\subsection{Resolution limit}
A mode cannot be tracked once its lobe amplitude falls below roughly one cell
radius, i.e. below $a_k \sim a/R$. In that regime the measured decay flattens and
\emph{understates} the true rate. Any dispersion measurement must therefore keep
the seeded amplitude well above $a/R$ throughout the fit window; this is enforced
as a validity condition in \texttt{experiments/giverso\_dispersion.py}, alongside
checks that the front is advancing, that the colony is not shedding, that the
quasi-steady nutrient field reached a fixed point, and that the advective CFL
number stays below one.

\subsection{Effective capillary length}
For a growing front stabilized by surface tension, the marginal mode $k_0$ (where
$\lambda$ changes sign) scales as $k_0\sim\sqrt{R/d_0}$ with $d_0$ the capillary
length. Measuring $k_0$ therefore yields
\begin{equation}
  d_0 \;\sim\; R/k_0^2,
\end{equation}
an order-of-magnitude estimate of the model's effective surface tension expressed
as a length. In CellFlow's soft-particle mechanics this comes out at a few cell
radii (Section~\ref{sec:limits}), which bounds how many fingers a colony of a
given size can support.

\subsection{Emergent surface tension: energy versus stress}\label{sec:kb}

Whether the model's \emph{own} mechanics carries a surface tension cannot be
settled by watching a shape relax, because a shape that does not relax may have
no surface tension or may have one and be unable to reach the lower-energy state.
The mechanical route separates them. For a flat interface with normal along $x$,
\begin{equation}
  \gamma = \int \left[\sigma_{xx}(x) - \sigma_{yy}(x)\right]\mathrm{d}x ,
\end{equation}
(Kirkwood--Buff), with the stress taken as the pair virial
$\sigma_{ab}=-A^{-1}\sum_{i<j} f_{ij,a} r_{ij,b}$ binned along $x$. This is read
off the current configuration and so is immune to kinetic arrest.

Measured on a JKR slab, the result explains the shape behaviour:
\begin{center}
\begin{tabular}{lrr}
\toprule
pack & $\gamma$ (stress) & $\gamma$ (energy, bond counting) \\ \midrule
perfect lattice at equilibrium & $0.00000$ & $0.041$ \\
jittered + relaxed, $w=0.3$    & $-0.020$  & $0.041$ \\
jittered + relaxed, $w=1.0$    & $-0.023$  & $0.328$ \\
\bottomrule
\end{tabular}
\end{center}

A pair potential on a lattice at its own equilibrium spacing has \emph{every}
bond at zero force: in a hex lattice the bulk lattice constant \emph{is} the pair
equilibrium, since all six neighbours are equidistant and $6F(s)=0$ forces
$F(s)=0$. A surface cell has fewer neighbours, but each remaining one is still at
that distance, so it too is force-free. The interface therefore carries positive
surface \emph{energy} (bonds are missing) and zero surface \emph{stress}. These
are distinct quantities related by Shuttleworth,
$f=\gamma+\mathrm{d}\gamma/\mathrm{d}\varepsilon$: for a liquid they are equal and
the interface pulls itself in; for a solid they differ, and only the stress drives
shape change. The tissue as modelled is a solid.

The practical consequence is that an \emph{emergent} surface tension requires a
change of rheology---neighbour exchange, cell deformability, or persistent active
motility---not a change of force law, and that $\gamma_{\text{stress}}$ rising to
meet $\gamma_{\text{energy}}$ is the criterion for having achieved it.

\section{Verification and validation}\label{sec:vv}

CellFlow ships an automated test suite (266 tests) including:
\begin{itemize}
  \item \textbf{Solver verification:} a Method of Manufactured Solutions test
        recovers a multi-mode divergence-free field to machine precision for
        both constant and spatially-varying drag.
  \item \textbf{Growth-driven expansion} (Section~\ref{sec:growthsource}): a
        manufactured potential is recovered exactly; the solved divergence equals
        the mean-removed source (to machine precision for a band-limited source,
        and otherwise to that source's Nyquist content, which is pinned to its
        analytic value rather than absorbed by a loose tolerance); the flux
        through a circle equals the enclosed source (Gauss); and the far field
        follows the 2D point-source law $u_r = Q/2\pi r$ to within $3\%$. The
        source contribution superposes linearly and is independent of $\mu$ and
        $\alpha$, as the pressure-absorption argument requires.
  \item \textbf{Front analysis} (Section~\ref{sec:front}): mode amplitudes are
        exact on analytic fronts; growth rates are recovered from synthetic
        colonies with a known $\lambda$; and a neutral front polluted with a
        growing number of detached cells returns $\lambda\approx 0$, which is the
        specific artifact the connected-cluster filter exists to prevent.
  \item \textbf{Convergence:} grid refinement studies confirm that the IBM
        self-mobility and hydrodynamic interactions converge; the timestep error
        scales as $\dt^1$ (forward Euler); the advection scheme conserves mass
        in the interior.
  \item \textbf{Transport solvers:} the implicit ADI diffusion is verified
        unconditionally stable at $\dt$ far beyond the explicit CFL limit, exact
        on the uniform steady state, mass-conserving under Neumann, and in
        agreement with the explicit scheme over a transient at small $\dt$.
        Multi-timescale: sub-stepping keeps explicit diffusion bounded at $4\times$
        the CFL limit; the coarse fluid cadence reproduces the every-step run on
        step~0 and remains stable with the cached field.
  \item \textbf{Physics checks:} screening monotonicity; mechanotransduction
        alignment to the correct strain axis; cell-shape response to uniaxial vs
        isotropic stress; area-conserving division; Michaelis--Menten uptake
        (reduces uptake at high concentration, recovers the linear law as
        $K_m\to\infty$, saturates, stays non-negative).
  \item \textbf{Reproducibility:} with a seed, runs are bit-for-bit identical
        and independent of the thread count.
\end{itemize}
\emph{Verification} (solving the equations correctly) is extensive;
\emph{validation} against experimental data is the user's responsibility for any
quantitative study, and physical units must be assigned by the user.

\section{Limitations}\label{sec:limits}

\begin{itemize}
  \item Two spatial dimensions only. Boundaries are either periodic with
        screening as a confinement surrogate, or free-slip walls on all four
        sides (Section~\ref{sec:freeslip}); explicit no-slip walls are not yet
        available.
  \item Overlap resolution is a geometric projection independent of viscosity;
        dense-tissue mechanics are therefore not viscosity-resolved.
  \item The differential-adhesion law is qualitative (not calibrated to a
        measured surface tension).
  \item The poroelastic Picard solver is efficient only for moderate drag
        contrast.
  \item Results are in arbitrary units unless the user assigns a unit system.
\end{itemize}

\paragraph{Interfacial patterning: what actually bounds it.}
This is worth stating explicitly because it constrains a whole class of studies,
and because an earlier version of this section was wrong in a way that mattered.

\emph{Withdrawn.} Previous editions gave a capillary length $d_0\sim R/k_0^2$ of
$3$--$8$ cell radii, a scaling $k^\ast\sim\sqrt{R/d_0}$, and a resulting estimate
of $10^5$--$10^7$ cells to resolve a branched morphology. The scaling law omitted
the nutrient diffusion length; measured over four sweeps it is
$k_0 \propto R/\sqrt{a\ell}$, the geometric mean of cell radius and diffusion
length, which lowered the cell-count estimate by about an order of magnitude. The
$10^5$--$10^7$ figure should not be cited.

\emph{What replaced it.} The binding constraint is not band width but the
\emph{shape} of the dispersion relation. Varying only the expansion rate---basal
metabolism slows net growth while leaving the nutrient field untouched---gives
both coefficients scaling together, with their ratios constant to $0.7\%$:
\begin{equation}
  \lambda(k) \;\approx\; \frac{\dot R}{R}\,\bigl[\,10 - 0.6\,k\,\bigr] .
\end{equation}
The marginal mode $k_0=\lambda_0/c\approx 16.7$ is therefore a \emph{pure
number}, independent of growth rate. Mullins--Sekerka is
$\lambda = Vk - \Gamma k^3$, in which the front velocity multiplies only the
destabilising term while $\Gamma$ is a material property; that asymmetry is what
selects a wavelength, $k^\ast=\sqrt{V/3\Gamma}$. The model as originally
formulated had no growth-rate-independent term at all, which is why four
independent parameters---nutrient level, growth rate, colony size and substrate
drag---each change the rate and leave the shape alone.

\emph{Current status.} The capillary term now exists
(Section~\ref{sec:surften}): raising $\sigma$ collapses $k_0$ from $14.4$ to
$4.2$ and roughly doubles the Mullins--Sekerka fit quality, confirming a genuine
$-\Gamma k^3$ cutoff. It moves $k^\ast$ \emph{down}, however, as
$k^\ast=\sqrt{V/3\Gamma}$ requires. The outstanding gap is therefore a
destabilising term that \emph{grows} with $k$; the model's is $k$-independent
(the geometric instability of any expanding circle). Note also that the fluid
velocity law suppresses force response by $1+(k\delta)^2$
(Section~\ref{sec:motion}), which acts directly against such a term.

See \texttt{docs/giverso\_replication.md} for the measurements and their
retraction history.

\end{document}
